\section*{第 \ref{ch:g11:vect} 章——平面中的向量与直线}

\begin{solution}{exo:g11:vect:1}
$\vect{AB}\,(3, 4)$, $\vect{AC}\,(-3, 2)$,
$\vect{AB} + \vect{AC}\,(0, 6)$,
$2\vect{AB} - 3\vect{AC}\,(6 + 9,\ 8 - 6) = (15, 2)$。线段
$[BC]$ 的\omterm{prop:g10:coordgeom:midpoint}{中点}为 $\left(\frac{5 + (-1)}{2}, \frac{3+1}{2}\right) = (2, 2)$。
\end{solution}

\begin{solution}{exo:g11:vect:2}
$\det = 4 \times 3 - (-2) \times (-6) = 12 - 12 = 0$：\omterm{def:g10:vectors:collinear}{共线}（事实上
$\vec v = -\frac12 \vec u$）。

$\det = 2 \times 7 - 3 \times 5 = -1 \neq 0$：不\omterm{def:g10:vectors:collinear}{共线}。

$\det = 0 \times (-8) - 0 \times 3 = 0$：\omterm{def:g10:vectors:collinear}{共线}（两\omterm{def:g11:vect:vector}{向量}均竖直）。
\end{solution}

\begin{solution}{exo:g11:vect:3}
过点 $A(2,1)$，方向为 $(3, -1)$：由 $(-b, a) = (3, -1)$ 得
$a = -1$，$b = -3$，于是 $-x - 3y + c = 0$；代入 $A$ 得 $-2 - 3 + c =
0$，$c = 5$。两边乘以 $-1$，得 $x + 3y - 5 = 0$。

过 $B(0,4)$ 与 $C(2,0)$：$\vect{BC}\,(2, -4)$，于是 $a = -4$，
$b = -2$：$-4x - 2y + c = 0$；代入 $B$ 得 $-8 + c = 0$，$c = 8$。
两边除以 $-2$ 化简，得 $2x + y - 4 = 0$。代入 $C$ 检验：$4 + 0 - 4 = 0$。
\end{solution}

\begin{solution}{exo:g11:vect:4}
\omterm{def:g11:vect:direction}{方向向量}为 $(-b, a) = (2, 3)$。与坐标轴的交点：令 $x = 0$，得
$y = 3$，得点 $(0, 3)$；令 $y = 0$，得 $x = -2$，点为 $(-2, 0)$。
对 $P(2,6)$：$3 \times 2 - 2 \times 6 + 6 = 0$，故 $P \in d$。对 $Q(1,4)$：
$3 - 8 + 6 = 1 \neq 0$，故 $Q \notin d$。
\end{solution}

\begin{solution}{exo:g11:vect:5}
$\vect{AB}\,(3, 6)$ 与 $\vect{AC}\,(-3, -6)$：
$\det = 3 \times (-6) - (-3) \times 6 = 0$，故 $A$、$B$、$C$ \omterm{def:g10:vectors:collinear}{共线}。

$\vect{DE}\,(2, 3)$ 与 $\vect{DF}\,(5, 8)$：
$\det = 2 \times 8 - 5 \times 3 = 1 \neq 0$：$D$、$E$、$F$ 不
\omterm{def:g10:vectors:collinear}{共线}。
\end{solution}

\begin{solution}{exo:g11:vect:6}
第一组：两式相加，$(2x + y - 7) + (x - y + 1) =
3x - 6 = 0$，故 $x = 2$，于是 $y = x + 1 = 3$，两直线相交于 $(2, 3)$
（检验：$4 + 3 - 7 = 0$）。

第二组：$ab' - a'b = 1 \times 4 - (-2) \times (-2) = 0$，故
两直线平行；将 $e$ 乘以 $-2$ 得 $-2x + 4y - 6 = 0$，
它与 $e'$ 的常数项不同（$-6 \neq 1$），
故二者平行且不重合，无交点。
\end{solution}

\begin{solution}{exo:g11:vect:7}
对 $d: 2x - y + 3 = 0$，\omterm{def:g11:vect:direction}{方向向量}为 $(-b, a) = (1, 2)$，点
$(0, 3)$ 在 $d$ 上，故 $d = \{(t,\ 3 + 2t) : t \in \R\}$。

对过点 $(1, -3)$、\omterm{def:g11:vect:direction}{方向向量}为 $(2, 1)$ 的直线：
$a = 1$，$b = -2$，于是 $x - 2y + c = 0$；代入 $(1, -3)$：
$1 + 6 + c = 0$，$c = -7$。\omterm{def:g10:algebra:equation}{方程}为 $x - 2y - 7 = 0$（检验 $t = 1$，
点 $(3, -2)$：$3 + 4 - 7 = 0$）。
\end{solution}

\begin{solution}{exo:g11:vect:8}
平行直线的 $(a, b)$ 相同，设其\omterm{def:g10:algebra:equation}{方程}为 $4x - 3y + c = 0$；过
点 $P(1, -2)$：$4 + 6 + c = 0$，故 $c = -10$，所求直线为
$4x - 3y - 10 = 0$。

要使 $mx + 2y - 5 = 0$ 与 $d$ 平行，需
$ab' - a'b = 4 \times 2 - m \times (-3) = 8 + 3m = 0$，故
$m = -\frac83$。
\end{solution}

\begin{solution}{exo:g11:vect:9}
线段 $[AB]$ 的\omterm{prop:g10:coordgeom:midpoint}{中点}为 $M\!\left(\frac{-1+3}{2}, \frac{0+2}{2}\right)
= (1, 1)$。中线为直线 $(CM)$，其中 $C(1, -4)$；两点
的\omterm{def:g10:coordgeom:system}{横坐标}均为 $1$，故中线为竖直线
\[
x - 1 = 0 .
\]
\end{solution}

\begin{solution}{exo:g11:vect:10}
$\det(\vec u, \vec v) = m(m - 1) - 3 \times 2 = m^2 - m - 6$。\omterm{def:g11:vect:det}{行列式}为零
当 $m^2 - m - 6 = 0$ 时：$\Delta = 1 + 24 = 25$，根为
$m = \frac{1 \pm 5}{2}$，即 $m = 3$ 或 $m = -2$。
\end{solution}

\begin{solution}{exo:g11:vect:11}
\emph{1.} $I = \left(\frac12, 0\right)$。于是
$\vect{DI} = \left(\frac12 - 0,\ 0 - 1\right) = \left(\frac12, -1\right)$
且 $\vect{DJ} = \frac23\vect{DI} = \left(\frac13, -\frac23\right)$，故
\[
J = D + \vect{DJ} = \left(\frac13,\ \frac13\right).
\]

\emph{2.} $\vect{AJ}\left(\frac13, \frac13\right)$ 与
$\vect{AC}\,(1, 1)$：
$\det = \frac13 \times 1 - 1 \times \frac13 = 0$，故 $A$、$J$、$C$ 三点
\omterm{def:g10:vectors:collinear}{共线}——事实上 $\vect{AJ} = \frac13\vect{AC}$，即直线 $(DI)$
与对角线 $(AC)$ 的交点恰在距 $A$ 为其长度三分之一处。
\end{solution}

\begin{solution}{pb:g11:vect:1}
\textbf{1.} $x = 1 + 3t$，$y = 2 - t$。由第二式，
$t = 2 - y$；代入得 $x = 1 + 3(2 - y) = 7 - 3y$：
直角坐标\omterm{def:g10:algebra:equation}{方程}为 $x + 3y - 7 = 0$。

\textbf{2.} 对于 $ax + by + c = 0$，\omterm{def:g11:vect:direction}{方向向量}为
$(-b, a)$。此处为 $(5, 2)$。取一点：$y = 1$ 时 $x = 1$：
$(1, 1)$。参数\omterm{def:g10:algebra:equation}{方程}为 $x = 1 + 5t$，$y = 1 + 2t$。

\textbf{3.} 将 $x = 7 - 3y$ 代入 $2x - 5y + 3 = 0$：
$14 - 6y - 5y + 3 = 0$，故 $y = \frac{17}{11}$，且
$x = 7 - \frac{51}{11} = \frac{26}{11}$，所求点为
$\left(\frac{26}{11}, \frac{17}{11}\right)$。

\textbf{4.} \omterm{def:g11:vect:direction}{方向向量}分别为 $(5, 2)$ 与 $(3, -1)$：
$\det = 5 \times (-1) - 2 \times 3 = -11 \neq 0$，两直线相交，
与前面结论一致。对于第二组，\omterm{def:g11:vect:direction}{方向向量} $(5, 2)$ 与
$(10, 4)$ 的\omterm{def:g11:vect:det}{行列式} $\det = 0$，二者平行，又将
$2x - 5y + 3 = 0$ 加倍得 $-4x + 10y - 6 = 0 \neq -4x + 10y + 1
= 0$：两直线严格平行。

\textbf{5.} $\vect{AB}(3, 2)$，$\vect{AC}(9, 6)$：
$\det = 18 - 18 = 0$，三点\omterm{def:g10:vectors:collinear}{共线}。

\textbf{6.} 对 $P$：由 $x = 2t$，$y = 3 + t$ 得 $y = 3 +
\frac x2$。对 $Q$：由 $x = 10 - t$ 得 $t = 10 - x$，故
$y = 2(10 - x) - 1 = 19 - 2x$。令两路径相交：
$3 + \frac x2 = 19 - 2x$，解得 $x = 6.4$，$y = 6.2$，两路径
相交于 $(6.4,\ 6.2)$。

\textbf{7.} $P$ 经过该点的时刻满足 $2t = 6.4$：$t = 3.2$ 小时。$Q$
经过该点的时刻满足 $10 - t = 6.4$：$t = 3.6$ 小时。两者相差二十四分钟，故无
碰撞。注意：图表只能展示\emph{路径}；只有
参数\omterm{def:g10:algebra:equation}{方程}中的时间变量才能判断两个运动是否
在同一时刻占据同一点。

\textbf{8.} $D^2(t) = (3t - 10)^2 + (4 - t)^2 =
10t^2 - 68t + 116$：顶点处 $t = \frac{68}{20} = 3.4$ 小时，
此时 $D^2 = 0.4$，最近距离为
$D = \sqrt{0.4} \approx 0.63$ 英里，对应 $t = 3.4$。

\textbf{9.} 条件被违反：$0.63 < 1$。$D^2(t) < 1$ 即
$10t^2 - 68t + 115 < 0$，其中 $\Delta = 4624 - 4600 = 24$，根为
$\frac{68 \pm \sqrt{24}}{20} = 3.4 \pm 0.245$，即两船
在 $t \approx 3.16$ 到 $t \approx 3.64$ 之间过于接近——约
$29$ 分钟的违规时长。两船之中须有一船改变航向。

\textbf{10.} 每艘船的每个坐标都是 $t$ 的\emph{一次}
\omterm{def:g11:func:function}{函数}，故每个坐标之差也是\omterm{def:g10:reffunc:affine}{一次函数}，
而 $D^2$——即两个\omterm{def:g10:reffunc:affine}{一次函数}的平方和——是一个
\omterm{def:g11:quad:quadratic}{二次函数}（二次项系数非负）。最近
距离因此总可由顶点读出。

\textbf{11.} 与 $P$ 在 $t = 3$ 处相遇，即位于 $(6, 6)$：
$(3a, 3b) = (6, 6)$，得 $(a, b) = (2, 2)$，速率为
$\sqrt{8} \approx 2.8$ 节。

\textbf{12.} $D(0, 1)$，$I\left(\frac12, 0\right)$：方向
$\left(\frac12, -1\right)$，即 $(1, -2)$，\omterm{def:g10:algebra:equation}{方程}为
$2x + y = 1$。

\textbf{13.} 直线 $(AC)$：$y = x$。代入得 $2x + x = 1$：
$x = \frac13$，交点为 $\left(\frac13, \frac13\right)$，恰为
自 $A$ 到 $C$ 的三分之一处。

\textbf{14.} $B(1, 0)$，$K\left(\frac12, 1\right)$：方向
$\left(-\frac12, 1\right)$，\omterm{def:g10:algebra:equation}{方程}为 $2x + y = 2$。联立
$y = x$，得 $x = \frac23$，交点为
$\left(\frac23, \frac23\right)$。两条塞瓦线在 $(AC)$ 上截出
两个三等分点。定理由此得证，所费计算不过
两次三条直线的演算。

\textbf{15.} 任一平行四边形都可以映射为单位正方形：
取 $A$ 为原点，并以 $\vect{AB}$、$\vect{AD}$ 作为两坐标轴的单位\omterm{def:g11:vect:vector}{向量}。
这一选取保持定理所讨论的一切
性质——\omterm{prop:g10:coordgeom:midpoint}{中点}、\omterm{def:g10:vectors:collinear}{共线}，以及沿给定直线的比值
（均可用\omterm{def:g11:vect:vector}{向量}与数量表示），因此在正方形中证明命题
即对所有平行四边形均成立。
（这一映射\emph{不}保持长度与角度：这些量会被
坐标架扭曲。）

\textbf{16.} $E\left(1, \frac12\right)$，直线 $(AE)$：
$y = \frac x2$。对角线 $(DB)$ 自 $(0,1)$ 至 $(1,0)$：
$y = 1 - x$。联立 $\frac x2 = 1 - x$，得 $x = \frac23$，
交点为 $\left(\frac23, \frac13\right)$——又是
一个三等分点，这次位于另一条对角线上。\omterm{prop:g10:coordgeom:midpoint}{中点}塞瓦线
总是偏爱三等分点。

\textbf{17.} 前两条：$x + y = 4$ 与 $2x - y = 2$ 联立得
$3x = 6$，得 $(2, 2)$。第三条直线：$2 - 4 = -2$，满足。三线
共点于 $(2, 2)$。

\textbf{18.} $mx - y + 2 - 3m = m(x - 3) + (2 - y) = 0$：要使
该\omterm{def:g10:algebra:equation}{方程}对\emph{每个} $m$ 均成立，取 $x = 3$，
$y = 2$ 即可——事实上每条直线 $d_m$ 都经过 $(3, 2)$。这组直线
构成过定点的直线束，每个\omterm{def:g10:reffunc:affine}{斜率} $m$ 对应其中一条。

\textbf{19.} 与 $y = 2x + 7$ 平行：\omterm{def:g10:reffunc:affine}{斜率} $m = 2$，所求直线为
$2x - y - 4 = 0$。过原点：$2 - 3m = 0$，
即 $m = \frac23$。

\textbf{20.} 直角坐标\omterm{def:g10:algebra:equation}{方程}：代入一点即可判断是否在直线上。
斜截式：\omterm{def:g10:reffunc:affine}{斜率}与截距一目了然，\omterm{def:g10:functions:graph}{图像}随即可得。
参数\omterm{def:g10:algebra:equation}{方程}：内置“时钟”，可判断碰撞、最近
距离与拦截。\omterm{def:g11:vect:det}{行列式}：一次交叉相乘，
即可判定平行。所选坐标架：平行四边形
化为单位正方形，而一个经典的三等分定理
化为三条直线的交点问题。
\end{solution}
